% Complete source Kolmogorov extension and measurable coding proofs.
\section{Countable products and real codings: proofs}
\label{sec:annex-p15-countable-product-measures}

The source is Durrett's fixed author version
\cite[pp.~53--54; app.~A.3, pp.~464--466]{DurrettPTEAuthor2019}.
Use the two cylinder algebras and all coding notation in
Chapter~\ref{sec:p15-countable-product-measures}.

\begin{annexitem}{annex:p15.cylinder-consistency}{thm:p15-cylinder-consistency}{The source algebra and coherent auxiliary Borel-cylinder masses}
\begin{proof}
Lift any two cylinders to a common prefix dimension by adding
unrestricted real coordinates. Rectangular cylinders are closed under
intersection, apart from a possible empty intersection, which is included.
The complement of a rectangular cylinder is a finite disjoint union
of rectangular cylinders: partition by the first coordinate that fails
its interval condition, and then by whether it lies to the left or
right of that interval. Earlier coordinates obey their original
conditions, later coordinates are unrestricted. Each one-dimensional
complement has at most two permitted interval pieces, including the
empty cases at infinite endpoints. The whole space is a rectangle.
Thus \(\mathcal S\) is a semialgebra, and its finite-disjoint-union
collection \(\mathcal A\) is an algebra by the earlier semialgebra
theorem.

The same common-prefix operation shows that \(\mathcal C\) is an
algebra: complements use \(B^c\), and finite unions use the union
of the lifted Borel bases. Clearly \(\mathcal A\subseteq\mathcal C\).
For fixed \(n\), the sets \(B\subseteq\mathbb R^n\) whose
inverse images under \(\pi_n\) belong to \(\mathcal F_\infty\)
form a sigma-algebra containing the half-open rectangles. Those
rectangles generate \(\mathcal B(\mathbb R^n)\): each bounded
open rectangle is an increasing union of half-open rectangles with
upper endpoints shifted inward, and rational open rectangles form
a Borel-generating basis. Conversely, half-open rectangles are Borel.
Thus every
Borel cylinder belongs to \(\mathcal F_\infty\).
Since \(\mathcal S\subseteq\mathcal A\subseteq\mathcal C\),
both generated sigma-algebras equal \(\mathcal F_\infty\).

The pushforward of \(\mu_{n+1}\) under its first-\(n\)
coordinate projection is a probability measure on \(\mathbb R^n\).
Under the bounded-rectangle consistency hypothesis it agrees with
\(\mu_n\) on the bounded half-open rectangle pi-system, which
generates the Borel sets. The probability uniqueness theorem
\ref{thm:p15-measure-uniqueness} gives equality on every Borel
set. Iteration proves
\(\mu_m(B\times\mathbb R^{m-n})=\mu_n(B)\) for every
\(m\geq n\), including the original unbounded rectangles.

Suppose \(\pi_n^{-1}(B)=\pi_m^{-1}(D)\), and lift both
bases to dimension \(r=\max(m,n)\). The lifted Borel sets
are equal: any point of \(\mathbb R^r\) extends to an infinite
sequence by appending zeros, so a difference of the lifted sets
would give a difference of the cylinders. Their \(\mu_r\)
values agree, and consistency identifies those values with
\(\mu_n(B)\) and \(\mu_m(D)\). Hence \(p_0\) is well defined.

For finitely many disjoint Borel cylinders, lift them to one prefix
dimension. Their bases are disjoint by the same surjectivity argument.
Finite additivity of the probability in that dimension gives finite
additivity of \(p_0\). Its empty value is zero and whole-space
value one. Decomposing a larger cylinder into a smaller one and
its difference gives monotonicity. Disjointifying a finite union
gives finite subadditivity. This establishes every use of \(p_0\)
on closed-base cylinders below without assuming countable additivity.
\end{proof}
\end{annexitem}

\begin{annexitem}{annex:p15.cylinder-compact-approximation}{thm:p15-cylinder-compact-approximation}{Bounded boxes, tightness and inner rectangle approximation}
\begin{proof}
A bounded sequence of real numbers has a convergent subsequence.
Indeed, repeatedly bisect a closed bounded interval containing it
and retain a half containing infinitely many terms. Choose successively
later terms in these nested intervals. The interval lengths tend to
zero, and completeness of the real order gives their common limit;
the chosen subsequence converges to it. In a closed bounded rectangle
in \(\mathbb R^d\), perform this extraction for each of the
finitely many coordinates. The resulting subsequence converges in
all coordinates, hence in Euclidean distance, and its limit remains
in the closed rectangle.

This sequential conclusion implies the stated open-cover compactness
here. Any open cover of the rectangle has a countable subcover:
for each rational open box lying inside some cover member and meeting
the rectangle, select one such member. Rational boxes form a countable
basis, so these selected members cover every point of the rectangle.
If a countable cover \(U_1,U_2,\ldots\) had no finite subcover,
choose \(x_n\) outside its first \(n\) members. A convergent
subsequence has a limit \(x\) in the rectangle. Some \(U_j\)
contains \(x\), and openness puts all sufficiently late terms of
that subsequence in \(U_j\), contradicting their exclusion once
their original indices exceed \(j\). Thus the cover has a finite
subcover. Empty or degenerate closed rectangles cause no exception.
Finite unions of compact sets are compact by taking a finite
subcover of each member. A closed subset of a compact set is
compact by adding its open complement to any cover. Closed
rectangles, and their finite unions, are closed by the coordinate
inequalities defining them.

The closed cubes \([-m,m]^d\), \(m\geq1\), increase to
\(\mathbb R^d\). Continuity from below of any Borel probability
gives masses increasing to one, so these compact cubes prove tightness.

For a nonempty rectangle \(R=\prod_{i=1}^d I(a_i,b_i)\), define
closed bounded approximating rectangles \(K_m\) as follows. In
coordinate \(i\), use lower endpoint
\(\max(-m,a_i+1/m)\) when \(a_i\) is finite, and \(-m\)
when \(a_i=-\infty\). Use upper endpoint \(\min(m,b_i)\)
when \(b_i\) is finite, and \(m\) when \(b_i=\infty\).
If a lower endpoint exceeds its upper endpoint, use the empty set.
Every \(K_m\) is contained in \(R\), the lower endpoints
decrease, and the upper endpoints increase. A point of \(R\)
eventually lies between all these bounds, so \(K_m\uparrow R\).
For a finite union \(B=\bigcup_{j=1}^r R_j\), take the union
of the respective approximations at the same \(m\). These are
finite unions of closed bounded rectangles, increase to \(B\),
and are contained in it. Finite measure subtraction and continuity
from below give \(\mu(B\setminus K_m)\to0\). Choose a
sufficiently large \(m\) for the requested error. An empty
original union is handled by \(K=\varnothing\). This proves
exactly the source's inner approximation for rectangle-algebra bases.
\end{proof}
\end{annexitem}

\begin{annexitem}{annex:p15.cylinder-continuity}{thm:p15-cylinder-continuity}{Compact finite-dimensional bases and a fully indexed diagonal limit}
\begin{proof}
The sequence \(p_0(B_n)\) decreases in \([0,1]\), so it has
a limit. Suppose that limit is \(\delta>0\). Write
\(B_n=\pi_{m_n}^{-1}(B_n^*)\), where \(B_n^*\) is a finite
union of half-open rectangles. By increasing the prefix dimensions
and adjoining unrestricted coordinates, choose \(m_n\) strictly
increasing with \(m_n\geq n\). This retains the original sets
and masses and avoids assuming that the original \(n\)-th set
already depends on only \(n\) coordinates.

By the compact approximation just proved, choose a finite union
of closed bounded rectangles \(K_n\subseteq B_n^*\) with
\[
 \mu_{m_n}(B_n^*\setminus K_n)<\delta\,2^{-(n+1)}.
\]
Set \(C_n=\pi_{m_n}^{-1}(K_n)\) and
\(D_n=\bigcap_{j=1}^n C_j\). These are in the Borel-cylinder
algebra \(\mathcal C\); they need not be in the initial
half-open rectangle algebra. Their \(p_0\) values are already
well defined and finitely additive by the first lemma.
Since \(B_n\subseteq B_j\) for \(j\leq n\),
\[
 B_n\setminus D_n\subseteq
                      \bigcup_{j=1}^n(B_j\setminus C_j).
\]
Finite subadditivity and the approximation errors give
\[
 p_0(B_n\setminus D_n)<\delta/2,
 \qquad p_0(D_n)>\delta/2.
\]
For the second inequality use \(D_n\subseteq B_n\), finite
additivity and \(p_0(B_n)\geq\delta\). Thus every \(D_n\)
is nonempty, and the \(D_n\) decrease.

Their finite-dimensional bases are
\[
 D_n^*=K_n\cap\bigcap_{j<n}
                    (K_j\times\mathbb R^{m_n-m_j}),
 \qquad D_n=\pi_{m_n}^{-1}(D_n^*).
\]
Each lifted earlier \(K_j\) is closed, being a projection
inverse image of a closed set. Thus \(D_n^*\) is a closed
subset of compact \(K_n\), and is compact and bounded. Choose
\(\omega^{(n)}\in D_n\). For each fixed coordinate \(i\),
all terms with \(n\geq i\) lie in \(D_i\); because
\(m_i\geq i\), their \(i\)-th coordinates lie in a bounded
projection of \(D_i^*\).

Extract an increasing subsequence of indices \(r_1(j)\) along
which coordinate one converges. Inductively choose an increasing
subsequence \(r_i(j)\) of \(r_{i-1}(j)\) along which
coordinate \(i\) also converges; discard finitely many indices
if needed so \(r_i(j)\geq\max(i,j)\). The real subsequence
lemma applies because that coordinate is eventually bounded.
Let its limit be \(\theta_i\). Define diagonal indices
\(s_k=r_k(k)\). They strictly increase: writing the next
subsequence as \(r_{k+1}(j)=r_k(t_j)\) with increasing
\(t_j\geq j\) gives \(s_{k+1}=r_k(t_{k+1})>r_k(k)=s_k\).
For every fixed \(i\), the terms with \(k\geq i\) form
a subsequence of the \(i\)-th chosen subsequence. Hence
\(\omega^{(s_k)}_i\to\theta_i\).

Fix \(n\). For all sufficiently large \(k\), \(s_k\geq n\),
so \(\omega^{(s_k)}\in D_n\). Its first \(m_n\)
coordinates converge to \((\theta_1,\ldots,\theta_{m_n})\).
Closedness of \(D_n^*\) puts that limit in \(D_n^*\).
Therefore the real sequence \(\theta=(\theta_1,\theta_2,\ldots)\)
belongs to every \(D_n\), and hence to every \(B_n\),
contradicting their empty intersection. The supposed positive
limit is impossible, so \(p_0(B_n)\downarrow0\).
This is the source's compact/diagonal argument with all cylinder
domains and coordinate dimensions specified.
\end{proof}
\end{annexitem}

\begin{annexitem}{annex:p15.kolmogorov-countable}{thm:p15-kolmogorov-countable}{Countable additivity followed by the measure extension}
\begin{proof}
Restrict \(p_0\) to \(\mathcal A\). Suppose pairwise disjoint
\(A_i\in\mathcal A\) have union \(A\in\mathcal A\).
The remainders \(B_n=A\setminus\bigcup_{i=1}^n A_i\)
belong to \(\mathcal A\), decrease to the empty set, and
finite additivity gives
\[
 p_0(A)=\sum_{i=1}^n p_0(A_i)+p_0(B_n).
\]
The continuity lemma makes the last term tend to zero. Taking
the limit of the nonnegative partial sums proves countable
additivity in every case whose union remains in the algebra.
Thus \(p_0|_{\mathcal A}\) is a premeasure with total mass
one. The full premeasure-extension theorem
\ref{thm:p15-premeasure-extension} gives a unique probability
measure \(P\) on \(\sigma(\mathcal A)=\mathcal F_\infty\)
extending that premeasure.

For fixed \(n\), \(\pi_n\) is measurable for this product
sigma-algebra, as proved in the first cylinder lemma. Its
pushforward under \(P\) is a probability on \(\mathbb R^n\)
agreeing with \(\mu_n\) on every half-open rectangle, since
those cylinders are in \(\mathcal A\). Uniqueness on the
generating rectangle pi-system proves agreement on every Borel
set \(B\). Hence \(P(\pi_n^{-1}(B))=\mu_n(B)\), including
all the auxiliary Borel cylinders used in the proof.

Any other probability with the prescribed rectangular values
agrees on their finite disjoint unions, and is another extension
of the same finite premeasure. It is therefore \(P\).
Each coordinate map is measurable because a coordinate inverse
image can be expressed using a sufficiently long prefix projection.
For any finite set of coordinate indices, choose a prefix length
at least their maximum and project its joint law \(\mu_m\)
to those indices. The consistency lemma makes this result
independent of the chosen larger prefix. This proves every
statement with the precise source index and sigma-algebra regime.
\end{proof}
\end{annexitem}

\begin{annexitem}{annex:p15.independent-sequence-existence}{thm:p15-independent-sequence-existence}{Prescribed marginal laws and the existing dyadic construction}
\begin{proof}
The finite products \(\mu_n=\eta_1\times\cdots\times\eta_n\)
exist by the finite product theorem. On any prefix rectangle,
adding a whole-line last coordinate multiplies its probability
by \(\eta_{n+1}(\mathbb R)=1\). Thus these probabilities
are consistent, and the proved countable extension applies.
Its coordinate marginal laws are \(\eta_i\). For any finite
subfamily, choose a prefix containing it and take the whole line
at unused coordinates. The prefix product law then gives exactly
the defining independence factorization for that subfamily.
This is mutual independence of the countable family.

For prescribed cumulative functions \(F_i\), the earlier
Stieltjes construction supplies \(\eta_i=\mu_{F_i}\).
The resulting rectangular cylinder probabilities are
\[
 P(a_i<X_i\leq b_i\text{ for }1\leq i\leq n)
                   =\prod_{i=1}^n(F_i(b_i)-F_i(a_i)),
\]
with the extensions \(F_i(-\infty)=0\), \(F_i(\infty)=1\)
at infinite endpoints. This is Example A.3.2.

The events \(A_n=\bigcap_{i=1}^n\{X_i\in E_i\}\)
decrease to the stated infinite-coordinate event. Their probabilities
are the finite products by independence. Continuity from above
of the now constructed probability measure gives the displayed
decreasing-product limit.

Finally, the map from the original uniform space to real sequences
whose coordinates are its dyadic bits is measurable: the inverse
image of each rectangular cylinder is a finite intersection of
measurable bit events. Its law on \(\mathcal F_\infty\)
has the fair binary product finite-dimensional distributions
already proved in Theorem~\ref{thm:p15-dyadic-independent-coins}.
Uniqueness of the extension identifies it with the canonical
measure just constructed. This is a comparison of two actual
probability laws after proving existence, not a substitute for
the countable-additivity argument.
\end{proof}
\end{annexitem}

\begin{annexitem}{annex:p15.sequence-borel-coding}{thm:p15-sequence-borel-coding}{The product metric and the exact interleaved binary image}
\begin{proof}
The series defining \(\delta\) is bounded by \(\sum_j2^{-j}=1\).
It is symmetric and nonnegative, and a zero value forces each
coordinate difference to be zero. Summing the real triangle
inequalities proves its triangle inequality. Thus it is a metric.
Each coordinate projection is continuous, because
\(|x_j-y_j|\leq2^j\delta(x,y)\).
Conversely, for a \(\delta\)-ball of radius \(\varepsilon\)
about \(x\), choose \(N\) with \(\sum_{j>N}2^{-j}<\varepsilon/2\)
and constrain \(|y_j-x_j|<\varepsilon/2\) for \(j\leq N\).
The resulting product neighborhood lies in that ball. These two
facts show equality of the metric and coordinate product topologies.
The product has a countable basis of finite-coordinate rational
interval conditions. Every open set is a countable union of such
basis sets, so the Borel sigma-algebra equals the coordinate
product sigma-algebra.

For \(t\in[0,1)\), the floor recurrence from the dyadic
construction, now including \(t=0\), gives binary digits
\(b_k(t)\in\{0,1\}\) and
\[
 t=\sum_{k\geq1}2^{-k}b_k(t),\qquad
 0\leq t-2^{-n}\lfloor2^nt\rfloor<2^{-n}.
\]
The digits are not eventually all one. If all digits after
position \(n\) were one, the first equality would give
\(t=2^{-n}(\lfloor2^nt\rfloor+1)\), contradicting that
floor. Thus they contain infinitely many zeros.
Conversely, a binary sequence with infinitely many zeros has
sum \(t\in[0,1)\), and its remainder after each \(n\)
digits is strictly less than \(2^{-n}\). Taking floors
therefore recovers exactly those digits. This proves the stated
canonical convention, including terminating dyadic expansions.

The formula \(\theta(j,k)\) enumerates the diagonals
\(j+k-1=1,2,\ldots\), with \(j\) increasing within
each diagonal. The diagonal of length \(r\) uses positions
\(r(r-1)/2+1\) through \(r(r+1)/2\). These finite
blocks are disjoint and cover all positive positions, so
\(\theta\) is a bijection. Every coordinate digit map is
Borel, its level sets being finite unions of half-open intervals.
Consequently the finite partial sums defining \(\phi_0\)
are measurable, and so is their nonnegative increasing limit.
The interleaved digit sequence has infinitely many zeros, since
each row does. Its value is less than one and its canonical
digits are exactly the interleaved digits.

The set \(E_0\) is Borel, being the displayed countable
intersection of countable unions of Borel digit level sets.
For \(x\in Q\), every row of its code consists of the
canonical digits of \(x_j\), so \(\phi_0(x)\in E_0\).
Its decoded coordinate is
\(D_j(\phi_0(x))=\sum_k2^{-k}b_k(x_j)=x_j\).
Thus \(\phi_0\) is injective and the displayed maps recover
all coordinates of its input.

For \(t\in E_0\), each decoded row has infinitely many
zeros. By the converse digit argument, \(D_j(t)\in[0,1)\)
and its canonical digits are exactly that row. Interleaving
these rows restores every canonical digit of \(t\), so
\(\phi_0(D(t))=t\). This proves surjectivity onto the exact
claimed image. Each \(D_j\) is measurable as an increasing
limit of finite Borel sums. The product measurability criterion
then makes \(D\) measurable into \(Q\). Zero codes the
all-zero input. One is outside the target interval, and any
code whose selected row is eventually all ones is explicitly
excluded from \(E_0\). Thus no inverse is being assigned
to a noncanonical or out-of-domain input.
\end{proof}
\end{annexitem}

\begin{annexitem}{annex:p15.metric-borel-representation}{thm:p15-metric-borel-representation}{The completed distance-coordinate construction and its Borel restriction}
\begin{proof}
If \(S=\varnothing\), the empty coding has all the claimed
properties. Suppose \(S\ne\varnothing\). A separable metric
space has a countable basis: balls with positive rational radii
centered in a countable dense set refine every open neighborhood.
Their intersections with \(S\) form a countable basis for
the subspace. Choose one point in each nonempty such intersection.
This gives a countable dense subset of \(S\), which can be
listed as a sequence \(q_1,q_2,\ldots\), repeating points
if the set is finite. This justifies the source's dense sequence
in \(S\), even when \(S\) is not closed or complete.

Let \(H=\overline S\) in \(M\). It is closed, so any
Cauchy sequence in \(H\) converges in \(M\) by completeness
and has its limit in \(H\). Thus \(H\) is complete, and
the same sequence \((q_j)\) is dense in \(H\).
For \(h(t)=t/(1+t)\), the function is increasing and
\[
 h(a)+h(b)-h(a+b)
   =\frac{ab(a+b+2)}{(1+a)(1+b)(1+a+b)}\geq0
                    \qquad(a,b\geq0).
\]
This inequality and the original triangle inequality show that
\(d=h\circ\rho\) is a metric, and \(d<1\).
For every \(\varepsilon>0\),
\(\rho(x,y)<\varepsilon\) exactly when
\(d(x,y)<\varepsilon/(1+\varepsilon)\).
This proves equality of topologies and equivalence of their
Cauchy and convergent sequences. In particular \(d\) is
complete on \(H\).

For \(x\in H\), the distance coordinates \(z_j=d(x,q_j)\)
belong to \([0,1)\). The triangle and reverse triangle
inequalities give
\( |z_i-z_j|\leq d(q_i,q_j)\leq z_i+z_j\),
and density gives \(\inf_jz_j=0\). Hence \(\psi(H)\subseteq C\).
Every coordinate of \(\psi\) is continuous; the preceding
product-metric argument makes \(\psi:H\to Q\) continuous.

Conversely, let \(z\in C\). For each \(k\geq1\), take
the least index \(n_k\) for which \(z_{n_k}<2^{-k}\),
which exists by the infimum condition. Then
\[
 d(q_{n_k},q_{n_l})\leq z_{n_k}+z_{n_l}<2^{-k}+2^{-l}.
\]
The selected points form a Cauchy sequence, so they converge
to some \(x\in H\). For fixed \(j\), the two coordinate
inequalities in \(C\) imply
\( |d(q_{n_k},q_j)-z_j|\leq z_{n_k}\).
Taking limits gives \(d(x,q_j)=z_j\), so \(\psi(x)=z\).
If two points \(x,y\) had these same coordinates, then
\(d(x,y)\leq d(x,q_j)+d(y,q_j)=2z_j\) for every \(j\).
The infimum condition makes \(d(x,y)=0\), hence \(x=y\).
This defines the inverse \(R(z)\) uniquely, with the stated
least-index Cauchy construction as an explicit description.

The inverse is continuous. Fix \(z=\psi(x)\) and
\(\varepsilon>0\). Choose \(j\) with \(z_j<\varepsilon/4\).
For \(z'=\psi(y)\in C\) with
\(|z'_j-z_j|<\varepsilon/2\),
\[
 d(x,y)\leq d(x,q_j)+d(y,q_j)=z_j+z'_j<\varepsilon.
\]
That coordinate condition is an open neighborhood in \(C\),
proving continuity at each \(z\). Thus \(\psi\) is
a homeomorphism onto \(C\).

The set \(C\) is Borel in \(Q\). Each of its coordinate
inequalities describes a closed set, and there are countably
many pairs \((i,j)\). Its last condition is
\[
 \{\inf_jz_j=0\}=\bigcap_{k\geq1}\bigcup_{j\geq1}
                                      \{z:z_j<1/k\},
\]
a Borel set because the coordinates are nonnegative. Since
\(S\) is Borel in \(M\) and is contained in the closed
subspace \(H\), it is Borel in \(H\). Therefore
\(K=\psi(S)=R^{-1}(S)\) is Borel in \(C\).
A Borel subset of a Borel subspace is Borel in the ambient
space: the trace formula expresses it as that subspace intersected
with an ambient Borel set. Hence \(K\) is Borel in \(Q\).

By the binary coding theorem, \(D:E_0\to Q\) is measurable
and \(E_0\) is Borel in \(\mathbb R\). Consequently
\(E=D^{-1}(K)\) is Borel in \(E_0\) and thus in
\(\mathbb R\). The composite \(\phi=\phi_0\circ\psi|_S\)
is a measurable bijection onto exactly this \(E\), because
each factor is a bijection onto its displayed image. Its inverse
is \(R\circ D\) restricted to \(E\). This map is measurable
into \(H\), takes values in \(S\), and is consequently
measurable for the trace Borel sigma-algebra of \(S\).
All image and inverse claims are now established. The use of
\(H=\overline S\) supplies the completeness needed for the
distance-coordinate image argument while retaining the source's
dense sequence in \(S\) and its actual interleaving construction.
\end{proof}
\end{annexitem}

\begin{annexitem}{annex:p15.coded-space-extension}{thm:p15-coded-space-extension}{Concentration on the coded coordinates and transport back}
\begin{proof}
The finite product of the maps \(\phi_i\) is a measurable
isomorphism from \(\prod_{i=1}^n S_i\) onto
\(E^{(n)}=\prod_{i=1}^nE_i\). This follows from the
coordinate measurability criterion, applied to both forward
and inverse maps. The product sigma-algebra on \(E^{(n)}\)
is the trace Borel sigma-algebra from \(\mathbb R^n\):
both are generated by intersections of Borel rectangles with
\(E^{(n)}\). The set \(E^{(n)}\) itself is Borel.

Push each \(\nu_n\) forward under this coordinate coding,
viewed as a map into \(\mathbb R^n\), obtaining \(\mu_n\).
The coding commutes with prefix projection. Thus the assumed
consistency of the \(\nu_n\) implies consistency of the
\(\mu_n\), and \(\mu_n(E^{(n)})=1\).
Kolmogorov's real-coordinate theorem supplies a probability
\(P\) on \((\mathbb R^{\mathbb N_+},\mathcal F_\infty)\)
with these finite-dimensional laws.

The set
\[
 \Omega_E=\prod_{i\geq1}E_i
          =\bigcap_{i\geq1}\{\omega:\omega_i\in E_i\}
\]
is measurable, since each \(E_i\) is Borel. Each event
in the intersection has probability one, by the concentration
of the corresponding finite-dimensional law. Their countable
intersection therefore has probability one. Restrict \(P\)
to \(\Omega_E\) with its trace sigma-algebra.
On that space define
\(T(\omega)=(\phi_i^{-1}(\omega_i))_{i\geq1}\).
Every coordinate is measurable; the product sigma-algebra on
\(\prod_i S_i\) is generated by those coordinates, so
\(T\) is measurable. Its pushforward is a probability measure
\(Q\) on the required space.

For a measurable first-\(n\) event \(B\), its image under
the finite coding is Borel in \(E^{(n)}\), because it is
the inverse image of \(B\) under the measurable inverse
coding. Since \(E^{(n)}\) is Borel in \(\mathbb R^n\),
that image is a Borel subset of \(\mathbb R^n\).
The construction of \(P\), followed by its concentration
on \(\Omega_E\), therefore gives
\(Q(\text{first }n\text{ coordinates lie in }B)=\nu_n(B)\).
This proves every required finite-dimensional distribution.

Two probabilities with those distributions agree on all finite
measurable cylinders. Those cylinders are a pi-system generating
the coordinate product sigma-algebra, so probability uniqueness
gives equality. Applying the preceding representation theorem
to the specified Borel subsets of complete separable metric
spaces supplies the final claimed class. The proof used Borel
images \(E_i\) to make the full-measure set \(\Omega_E\)
measurable; it did not presume this for arbitrary measurable spaces.
\end{proof}
\end{annexitem}
